\chapter{总结与展望}\label{chap:conclusion}

\section{总结}

本文围绕高校多校区办学背景下的校园出行问题，完成了一个基于 Spring Boot 的校内综合交通平台原型。论文所处理的并不是抽象意义上的“智慧交通”命题，而是高校内部较为具体的几类问题：班车信息长期分散，拼车协作缺少稳定流程，许多真实的出行行为也难以沉淀为后续可继续利用的数据。围绕这些问题，本文完成了需求分析、系统设计、功能实现和测试验证。

从系统组织方式看，平台以前后端分离作为基本实现思路，由微信小程序承担统一用户入口，后端按用户、拼车、论坛、聊天、文件和交通等业务拆分服务。这样处理并不是为了让架构看起来更复杂，而是因为不同业务的数据特点和处理方式差别明显。用户和拼车数据存在状态变化与一致性要求，论坛、聊天和时刻表数据则更强调结构灵活与读取便利。因此，系统在数据层采用了 MySQL 与 MongoDB 组合持久化的方案，让不同类型的数据能够放在更合适的存储环境中。

本文实现工作的重心放在拼车链路上。系统没有把拼车理解为单次信息发布，而是把它组织为订单发布、候选查询、申请审批、行程完成和互评反馈这一连续过程。反馈结果会继续写回用户信用统计，并在后续匹配中再次发挥作用。针对候选排序问题，本文实现了“规则筛选 + 随机森林排序”的二层匹配机制。规则层先排除不满足基本约束的订单，再由模型对高质量候选做进一步区分。这样安排以后，匹配结果既不脱离校园场景中的基本约束，也不会停留在简单平铺的层面。

论坛、聊天和交通模块虽然不处于拼车主链中心，但它们补足了经验交流、即时确认和基础查询等环节，使平台能够形成较完整的使用闭环。结合测试结果来看，在当前原型规模下，系统已经能够完成主要业务流程，接口交互与数据更新逻辑也保持了基本稳定。对本文而言，更重要的收获不在于单独提出某一种模型，而在于把原本分散的校园交通服务组织到同一平台中，并验证其具备实际运行的基本条件。
\section{展望}

本文完成了系统的主要设计与实现工作，但若从真实运行和后续研究的角度来看，平台仍有一些需要继续推进的地方。

比较直接的一点在于拼车匹配算法。当前模型训练主要基于模拟样本完成，足以支撑原型系统验证，却还不能完全代表真实校园出行行为。后续如果能够在系统运行过程中持续积累真实订单、取消记录、等待时长和用户反馈，训练样本会更接近实际场景，匹配质量也有望随之提高。

系统性能与部署层面也还有进一步完善的空间。现阶段原型系统已经完成基本服务拆分，但在服务治理和可观测性方面仍然偏轻量。随着访问量上升，统一网关、日志追踪、监控告警以及缓存机制的重要性都会明显增加。尤其是拼车查询这类包含多步计算和跨服务调用的接口，后续还可以通过缓存优化或异步处理来减轻高峰时段的响应压力。

平台如果面向更大规模用户持续运行，治理问题也会变得更具体。隐私保护、内容审核、异常行为识别和访问控制都不宜继续作为附加功能来看待，而会逐渐变成平台稳定运行的基础条件。当前系统已经具备基础身份认证和信用约束，但若进入长期运行阶段，仍需要补充更细的权限管理和审计机制。

总体来看，本文完成的平台已经具备较完整的基础框架，能够支撑校园综合交通服务中的主要流程。后续如果在真实数据积累、服务治理和平台运行机制上继续推进，其实际应用价值仍有进一步提升的空间。
